\chapter{Методы Монте-Карло}
\label{ch:mc}

\begin{quote}
\itshape
Чтобы найти ценность позиции, симулируйте множество партий из этой позиции и усредните результаты.
\end{quote}

\bigskip

Методы динамического программирования из главы~\ref{ch:dp} обладают большой мощью, однако требуют жёсткого предварительного условия: полного знания вероятностей переходов среды и функции вознаграждения.  Во многих практически важных задачах\,---\,карточных играх, управлении роботами, диалоговых системах\,---\,такая модель попросту недоступна или слишком велика для полного перебора всех состояний.  \emph{Методы Монте-Карло}\index{Monte Carlo methods} полностью устраняют это требование.  Вместо того чтобы вычислять функции ценности из математической спецификации среды, они \emph{обучаются} функциям ценности непосредственно по выборкам эпизодов взаимодействия.

Основная идея столь же стара, как и статистическое моделирование: если вы хотите узнать ожидаемый результат стохастического процесса, запустите его много раз и усредните результаты.  Применительно к обучению с подкреплением это означает генерацию полных эпизодов в рамках фиксированной стратегии, вычисление отдачи (дисконтированной суммы вознаграждений), следующей за каждым посещением состояния или пары состояние--действие, и усреднение этих отдач.  По мере роста числа эпизодов закон больших чисел гарантирует сходимость к истинной функции ценности.

Методы Монте-Карло занимают ключевое место в пейзаже обучения с подкреплением.  Они знаменуют переход от обучения на основе модели к \emph{безмодельному}\index{model-free} обучению и вводят фундаментальную дихотомию между \emph{обучением на текущей стратегии}\index{on-policy} и \emph{обучением на другой стратегии}\index{off-policy}, которая будет встречаться на протяжении всей остальной части книги.  Они требуют полных эпизодов, что делает их естественно пригодными для эпизодических задач; методы временной разности из главы~\ref{ch:td} позднее устранят и это ограничение.


% ===========================================================
\section{Предсказание методом Монте-Карло}
\label{sec:mc-prediction}
\index{Monte Carlo prediction}

\subsection*{Постановка задачи и обозначения}

На протяжении всей этой главы мы работаем с конечным эпизодическим МПР (см. главу~\ref{ch:mdp}) с пространством состояний $\cS$, пространством действий $\cA$, коэффициентом дисконтирования $\gamma \in [0,1)$ и фиксированной стратегией $\pi$.  Один эпизод представляет собой траекторию
\[
  S_0, A_0, R_1, S_1, A_1, R_2, \;\dots,\; S_{\tau-1}, A_{\tau-1}, R_{\tau}, S_{\tau} = \mathrm{T},
\]
где $\mathrm{T}$ --- терминальное состояние, а $\tau$ --- (случайная) длина эпизода.  \emph{Отдача}\index{return} начиная с шага $t$ равна
\begin{equation}
  \label{eq:mc-return}
  G_t = \sum_{k=0}^{\tau - t - 1} \gamma^k R_{t+1+k}.
\end{equation}
Функция ценности состояния задаётся как $v^{\pi}(s) = \E_\pi[G_t \mid S_t = s]$.

Мы предполагаем на протяжении всей главы, что: (i) эпизоды генерируются независимо в соответствии со стратегией $\pi$; (ii) каждый эпизод завершается почти наверное, $\Prob_\pi(\tau < \infty) = 1$; (iii) вознаграждения ограничены, $|R_t| \le R_{\max}$ п.н.; (iv) каждое интересующее нас состояние $s$ посещается с положительной вероятностью за один эпизод.

\subsection*{Монте-Карло по первому посещению и по каждому посещению}

В рамках одного эпизода состояние $s$ может посещаться несколько раз.  Два естественных способа усреднения соответствуют разным подходам к обработке повторных посещений.

\begin{definition}[Оценщик МК по первому посещению]
\label{def:first-visit-mc}
\index{first-visit Monte Carlo}
Оценщик \textbf{Монте-Карло по первому посещению} для $v^\pi(s)$ усредняет по эпизодам отдачу $G_t$, следующую за \emph{первым} шагом $t$, на котором $S_t = s$ в каждом эпизоде.  Если $G_1^{\mathrm{FV}}(s), G_2^{\mathrm{FV}}(s), \dots$ --- это отдачи, собранные из $n$ эпизодов, то оценщик равен
\begin{equation}
  \label{eq:fv-estimator}
  \widehat{v}_n^{\,\mathrm{FV}}(s) = \frac{1}{n} \sum_{i=1}^{n} G_i^{\mathrm{FV}}(s).
\end{equation}
\end{definition}

\begin{definition}[Оценщик МК по каждому посещению]
\label{def:every-visit-mc}
\index{every-visit Monte Carlo}
Оценщик \textbf{Монте-Карло по каждому посещению} усредняет отдачи, следующие за \emph{каждым} посещением $s$ по всем эпизодам.  Для эпизода $i$ пусть $K_i(s) = \sum_t \ind\{S_t^{(i)} = s\}$ --- число посещений, а $X_i(s) = \sum_t \ind\{S_t^{(i)} = s\} G_t^{(i)}$ --- суммарные отдачи.  Оценщик после $n$ эпизодов равен
\begin{equation}
  \label{eq:ev-estimator}
  \widehat{v}_n^{\,\mathrm{EV}}(s) = \frac{\sum_{i=1}^{n} X_i(s)}{\sum_{i=1}^{n} K_i(s)}.
\end{equation}
\end{definition}

Оба оценщика нацелены на одну и ту же величину $v^\pi(s)$.  Их различие носит статистический, а не определяющий характер.

\subsection*{Несмещённость МК по первому посещению}

\begin{theorem}[МК по первому посещению несмещён и состоятелен]
\label{thm:fv-unbiased}
При сделанных выше допущениях каждая отдача $G_i^{\mathrm{FV}}(s)$ является независимой одинаково распределённой случайной величиной со средним $v^\pi(s)$ и конечной дисперсией.  Следовательно:
\begin{enumerate}
  \item $\E\!\left[\widehat{v}_n^{\,\mathrm{FV}}(s)\right] = v^\pi(s)$ для всех $n \ge 1$ (несмещённость).
  \item $\widehat{v}_n^{\,\mathrm{FV}}(s) \xrightarrow[n \to \infty]{\mbox{п.н.}} v^\pi(s)$ (сильная состоятельность).
\end{enumerate}
\end{theorem}

\begin{proof}
\textbf{Независимость.}  Каждая $G_i^{\mathrm{FV}}(s)$ вычисляется по отдельному независимо сгенерированному эпизоду.  Следовательно, последовательность $G_1^{\mathrm{FV}}(s), G_2^{\mathrm{FV}}(s), \dots$ состоит из независимых случайных величин.

\textbf{Одинаковое распределение.}  Рассмотрим произвольный эпизод и пусть $T^*$ --- первый шаг, на котором $S_{T^*} = s$.  По свойству Маркова распределение траектории начиная с $T^*$ зависит только от текущего состояния $s$ и не зависит от истории $S_0, A_0, \dots, S_{T^*-1}$.  Поскольку $\pi$ фиксирована, отдача $G_{T^*}$ с этой точки имеет в точности распределение $G_t \mid S_t = s$ при стратегии $\pi$.  Таким образом, $G_i^{\mathrm{FV}}(s) \stackrel{d}{=} G_t \mid S_t = s$ для всех $i$.

\textbf{Конечная дисперсия.}  Поскольку $|R_t| \le R_{\max}$ и $\tau < \infty$ п.н.,
\[
  |G_t| \le R_{\max} \sum_{k=0}^{\tau-1} \gamma^k \le \frac{R_{\max}}{1-\gamma} < \infty \quad \mbox{п.н.}
\]
Следовательно, $G_i^{\mathrm{FV}}(s)$ ограничена п.н. и имеет конечную дисперсию $\sigma_s^2 < \infty$.

\textbf{Несмещённость.}  По линейности математического ожидания,
\[
  \E\!\left[\widehat{v}_n^{\,\mathrm{FV}}(s)\right]
  = \frac{1}{n} \sum_{i=1}^{n} \E\!\left[G_i^{\mathrm{FV}}(s)\right]
  = v^\pi(s).
\]

\textbf{Сильная состоятельность.}  Поскольку $G_1^{\mathrm{FV}}(s), G_2^{\mathrm{FV}}(s), \dots$ --- н.о.р. с конечным средним, усиленный закон больших чисел даёт
\[
  \widehat{v}_n^{\,\mathrm{FV}}(s) = \frac{1}{n}\sum_{i=1}^{n} G_i^{\mathrm{FV}}(s) \xrightarrow{\mbox{п.н.}} v^\pi(s). \qedhere
\]
\end{proof}

\begin{remark}
Дисперсия $\widehat{v}_n^{\,\mathrm{FV}}(s)$ равна $\sigma_s^2/n$ и убывает со стандартной скоростью Монте-Карло $1/\sqrt{n}$.  Это контрастирует, например, с бутстрэп-методами, которые могут достигать более высоких скоростей в гладких моделях.  Однако здесь не требуется никаких предположений о распределении\,---\,лишь ограниченность вознаграждений и эпизодическое завершение.
\end{remark}

\subsection*{Состоятельность МК по каждому посещению}

\begin{lemma}
\label{lem:ev-key-identity}
Для любого состояния $s$ и любого эпизода $\E[X_1(s)] = v^\pi(s)\,\E[K_1(s)]$.
\end{lemma}

\begin{proof}
По линейности и свойству башни,
\begin{align*}
  \E[X_1(s)]
  &= \sum_{t \ge 0} \E\!\left[\ind\{S_t = s\}\,G_t\right]
  = \sum_{t \ge 0} \E\!\left[\ind\{S_t = s\}\,\E[G_t \mid S_t]\right]\\
  &= v^\pi(s)\sum_{t \ge 0} \Prob(S_t = s)
  = v^\pi(s)\,\E[K_1(s)]. \qedhere
\end{align*}
\end{proof}

\begin{theorem}[МК по каждому посещению состоятелен]
\label{thm:ev-consistent}
Если $\E[K_1(s)] \in (0,\infty)$ и $\E[|X_1(s)|] < \infty$, то $\widehat{v}_n^{\,\mathrm{EV}}(s) \xrightarrow{\mbox{п.н.}} v^\pi(s)$.
\end{theorem}

\begin{proof}
Пары $(X_i(s), K_i(s))$ являются н.о.р. по эпизодам.  Усиленный закон больших чисел даёт
\[
  \frac{1}{n}\sum_{i=1}^n X_i(s) \xrightarrow{\mbox{п.н.}} \E[X_1(s)], \qquad
  \frac{1}{n}\sum_{i=1}^n K_i(s) \xrightarrow{\mbox{п.н.}} \E[K_1(s)] > 0.
\]
Знаменатель в конечном счёте положителен п.н., поэтому по теореме о непрерывном отображении (правило частных):
\[
  \widehat{v}_n^{\,\mathrm{EV}}(s) \xrightarrow{\mbox{п.н.}} \frac{\E[X_1(s)]}{\E[K_1(s)]} = v^\pi(s),
\]
где последнее равенство следует из леммы~\ref{lem:ev-key-identity}.
\end{proof}

\begin{remark}
МК по каждому посещению \emph{смещён} при любом конечном $n$, поскольку $\E[\sum X_i / \sum K_i] \ne \E[\sum X_i]/\E[\sum K_i]$ в общем случае (смещение оценщика-отношения).  Однако смещение исчезает при $n \to \infty$, поэтому МК по каждому посещению \emph{асимптотически} несмещён.  На практике МК по каждому посещению может иметь меньшую среднеквадратичную ошибку, чем МК по первому посещению, когда состояния часто повторяются внутри эпизода, поскольку использует больше данных.
\end{remark}

\subsection*{Инкрементная реализация}

Вычислять средние заново после каждого эпизода расточительно.  Вместо этого поддерживается текущий счётчик $N(s)$ и текущая оценка $V(s)$, которые обновляются после каждой наблюдаемой отдачи $G$:
\begin{equation}
  \label{eq:mc-incremental}
  N(s) \leftarrow N(s) + 1, \qquad
  V(s) \leftarrow V(s) + \frac{1}{N(s)}\bigl(G - V(s)\bigr).
\end{equation}
Это алгебраически эквивалентно скользящему выборочному среднему и требует $O(|\cS|)$ памяти.

\begin{algorithm}[ht]
\caption{Предсказание методом МК по первому посещению}
\label{alg:fv-mc-prediction}
\KwIn{Стратегия $\pi$, дисконт $\gamma$, число эпизодов $N_{\mathrm{ep}}$}
Инициализировать $V(s) \leftarrow 0$, $\mathit{Count}(s) \leftarrow 0$ для всех $s \in \cS$\;
\For{эпизод $= 1, 2, \dots, N_{\mathrm{ep}}$}{
  Сгенерировать эпизод $S_0, A_0, R_1, \dots, S_\tau$ по стратегии $\pi$\;
  $G \leftarrow 0$\;
  \For{$t = \tau - 1, \tau - 2, \dots, 0$ \upshape(в обратном порядке)}{
    $G \leftarrow \gamma G + R_{t+1}$\;
    \If{$S_t$ не встречается в $S_0, \dots, S_{t-1}$}{
      $\mathit{Count}(S_t) \leftarrow \mathit{Count}(S_t) + 1$\;
      $V(S_t) \leftarrow V(S_t) + \dfrac{1}{\mathit{Count}(S_t)}\bigl(G - V(S_t)\bigr)$\;
    }
  }
}
\Return $V$\;
\end{algorithm}

Обратный проход по эпизоду (строки 4--9) вычисляет все отдачи $G_t$ за время $O(\tau)$ с помощью рекуррентного соотношения $G_t = R_{t+1} + \gamma G_{t+1}$.

\subsection*{Резервная диаграмма МК}

Следующая диаграмма сравнивает резервные диаграммы динамического программирования и метода Монте-Карло.  В ДП одношаговое резервное обновление усредняет по всем возможным следующим состояниям с использованием модели.  В МК следуют по одной выборочной траектории до конца, до терминального состояния.

\begin{figure}[ht]
\centering
\begin{tikzpicture}[
    node distance = 0.7cm and 1.5cm,
    state/.style  = {circle, draw, minimum size=0.5cm, thick},
    action/.style = {circle, fill=black, minimum size=0.18cm, inner sep=0pt},
    terminal/.style = {rectangle, draw, minimum width=0.7cm, minimum height=0.35cm, thick, fill=gray!30},
    >=Stealth
  ]
  % --- Резервное копирование МК (слева) ---
  \node[state]   (s0)                       {$s$};
  \node[action]  (a0) [below=of s0]         {};
  \node[state]   (s1) [below=of a0]         {};
  \node[action]  (a1) [below=of s1]         {};
  \node[state]   (s2) [below=of a1]         {};
  \node[action]  (a2) [below=of s2]         {};
  \node[state]   (s3) [below=of a2]         {};
  \node[action]  (a3) [below=of s3]         {};
  \node[terminal](sT) [below=of a3]         {$\mathrm{T}$};

  \draw[->] (s0) -- (a0) node[midway, right=2pt] {\footnotesize $A_0$};
  \draw[->] (a0) -- (s1) node[midway, right=2pt] {\footnotesize $R_1$};
  \draw[->] (s1) -- (a1) node[midway, right=2pt] {\footnotesize $A_1$};
  \draw[->] (a1) -- (s2) node[midway, right=2pt] {\footnotesize $R_2$};
  \draw[->] (s2) -- (a2) node[midway, right=2pt] {\footnotesize $A_2$};
  \draw[->] (a2) -- (s3) node[midway, right=2pt] {\footnotesize $R_3$};
  \draw[->] (s3) -- (a3) node[midway, right=2pt] {\footnotesize $A_{\tau-1}$};
  \draw[->] (a3) -- (sT) node[midway, right=2pt] {\footnotesize $R_\tau$};

  \node[above=0.2cm of s0, font=\bfseries\small] {Резерв. МК};
  \node[below=0.15cm of sT, font=\small, align=center] {полный эпизод};

  % --- Резервное копирование ДП (справа) ---
  \node[state]   (ds0) [right=5cm of s0]      {$s$};
  \node[action]  (da1) [below left=0.7cm and 0.5cm of ds0]  {};
  \node[action]  (da2) [below=0.7cm of ds0]                  {};
  \node[action]  (da3) [below right=0.7cm and 0.5cm of ds0] {};
  \node[state]   (ds1) [below=0.65cm of da1]   {};
  \node[state]   (ds2) [below=0.65cm of da2]   {};
  \node[state]   (ds3) [below=0.65cm of da3]   {};

  \draw[->] (ds0) -- (da1);
  \draw[->] (ds0) -- (da2);
  \draw[->] (ds0) -- (da3);
  \draw[->] (da1) -- (ds1);
  \draw[->] (da2) -- (ds2);
  \draw[->] (da3) -- (ds3);

  \node[above=0.2cm of ds0, font=\bfseries\small] {Резерв. ДП};
  \node[below=0.15cm of ds2, font=\small, align=center] {один шаг, все ветви};

\end{tikzpicture}
\caption{Резервные диаграммы для метода Монте-Карло (слева) и динамического программирования (справа).  МК следует по одной выборочной траектории до терминального состояния; ДП разветвляется по всем возможным одношаговым переходам с использованием модели.}
\label{fig:mc-backup-diagram}
\end{figure}


% ===========================================================
\section{Разобранный пример: Блэкджек}
\label{sec:mc-blackjack}
\index{Blackjack}

Блэкджек (также называемый \emph{двадцать одно}) --- канонический пример, используемый для иллюстрации предсказания методом Монте-Карло в книге Саттона и Барто (2018).  Он особенно поучителен, поскольку полная модель среды (вероятности переходов) технически доступна, однако её трудно записать в явном виде, тогда как оценивание методом Монте-Карло легко реализовать путём симуляции карточных партий.

\subsection*{MDP блэкджека}

Игра ведётся против дилера.  Цель игрока --- набрать карты с суммой как можно ближе к 21, не превышая этого значения (не «перебрать»).  Картинки (валет, дама, король) считаются за 10; туз считается за 11, если при этом сумма не превысит 21, иначе туз считается за 1.

\paragraph{Состояние.}  Состояние --- тройка $(x, d, a) \in \{12,\dots,21\} \times \{1,\dots,10\} \times \{0,1\}$, где:
\begin{itemize}[leftmargin=*]
  \item $x$ = текущая сумма игрока (рассматриваются только суммы $\ge 12$; ниже 12 игрок всегда безопасно берёт карту);
  \item $d$ = открытая карта дилера (туз считается за 1);
  \item $a = 1$, если у игрока есть «рабочий» туз (туз, в данный момент считающийся за 11), $a = 0$ иначе.
\end{itemize}
Это даёт $|\cS| = 10 \times 10 \times 2 = 200$ нетерминальных состояний.

\paragraph{Действия.}  В каждом состоянии игрок выбирает \emph{стоп} (прекратить брать карты, $A=0$) или \emph{ещё} (взять ещё карту, $A=1$).

\paragraph{Вознаграждения.}  Эпизод завершается, когда игрок говорит «стоп» или перебирает.  После «стопа» дилер играет до конца по фиксированному правилу (берёт карты, пока сумма $< 17$).  Вознаграждение: $+1$ за победу игрока, $-1$ за проигрыш, $0$ за ничью.  Все промежуточные вознаграждения равны $0$.  Коэффициент дисконтирования $\gamma = 1$.

\paragraph{Стратегия.}  Оцениваем \emph{простую стратегию} $\pi$: говорить «стоп», если $x \ge 20$, иначе брать ещё.

\subsection*{Генерация эпизодов и вычисление отдачи}

\begin{example}[Пример эпизода блэкджека]
\label{ex:blackjack-episode}
\index{Blackjack!sample episode}
Рассмотрим следующий эпизод, сгенерированный в соответствии с простой стратегией:
\begin{center}
\small
\begin{tabular}{@{}clccc@{}}
\toprule
$t$ & Состояние & Действие & Сл.\,карта & $R_{t+1}$ \\
\midrule
0 & $(15, 6, 0)$ & Ещё  & 3  & 0 \\
1 & $(18, 6, 0)$ & Ещё  & 2  & 0 \\
2 & $(20, 6, 0)$ & Стоп & ---  & --- \\
\multicolumn{4}{l}{Дилер: $6 + 9 = 15$, берёт $4 \to 19$; игрок $20 > 19$.} & $+1$ \\
\bottomrule
\end{tabular}
\end{center}
Эпизод завершается при $\tau = 3$ и $R_3 = +1$.  Отдачи (при $\gamma=1$) равны:
\begin{align*}
  G_2 &= R_3 = +1, \\
  G_1 &= R_2 + G_2 = 0 + 1 = +1, \\
  G_0 &= R_1 + G_1 = 0 + 1 = +1.
\end{align*}
При МК по первому посещению этот эпизод вносит вклад:
\begin{itemize}[leftmargin=*]
  \item Отдача $G_0 = +1$ в оценку $v^\pi(15, 6, 0)$ (первое и единственное посещение).
  \item Отдача $G_1 = +1$ в $v^\pi(18, 6, 0)$.
  \item Отдача $G_2 = +1$ в $v^\pi(20, 6, 0)$.
\end{itemize}
\end{example}

Запуск $500{,}000$ таких эпизодов даёт точное представление о функции ценности.  Состояния с $x = 20$ или $x = 21$ имеют значения, близкие к $+0.6$ (игрок обычно выигрывает), тогда как состояния с $x \in \{12, \dots, 14\}$ против туза дилера имеют значения около $-0.3$.  В учебнике Саттона--Барто эти значения изображены в виде трёхмерных поверхностей, демонстрирующих интуитивно понятный ландшафт: высокие суммы против слабых карт дилера выгодны; низкие суммы против сильных карт дилера невыгодны.

Пример демонстрирует, почему метод Монте-Карло хорошо подходит для данной задачи: вычисление $p(s', r \mid s, a)$ в точности требует перечисления распределения карт, взятых из бесконечной колоды, тогда как симуляция раздачи тривиальна.


% ===========================================================
\section{Оценка функций ценности действий методом Монте-Карло}
\label{sec:mc-action-values}
\index{action-value function!Monte Carlo estimation}

Когда модель среды $p(s', r \mid s, a)$ доступна, функции ценности состояний $v^\pi$ достаточно для управления: зная $v^\pi$, можно вычислить жадную стратегию путём одношагового просмотра вперёд,
\[
  \pi'(s) \in \argmax_{a \in \cA} \sum_{s', r} p(s', r \mid s, a)\bigl[r + \gamma v^\pi(s')\bigr].
\]
Без модели такой просмотр невозможен.  Мы не можем вычислить правую часть, поскольку $p(s', r \mid s, a)$ неизвестна.  Поэтому в безмодельном управлении естественной целью является \emph{функция ценности действия}\index{action-value function}
\begin{equation}
  \label{eq:mc-q-def}
  Q^\pi(s, a) = \E_\pi[G_t \mid S_t = s,\; A_t = a].
\end{equation}
Зная $Q^\pi$, жадная стратегия просто $\pi'(s) = \argmax_{a} Q^\pi(s, a)$, и модель не нужна вовсе.

\paragraph{Оценка $Q^\pi$ методом Монте-Карло.}  Весь аппарат раздела~\ref{sec:mc-prediction} применяется дословно, с заменой состояний $s$ на пары состояние--действие $(s, a)$.  Оценщик МК по первому посещению для $Q^\pi(s, a)$ собирает отдачи с первого посещения $(s, a)$ в каждом эпизоде; оценщик по каждому посещению усредняет все такие отдачи.  Теоремы~\ref{thm:fv-unbiased} и~\ref{thm:ev-consistent} справедливы с идентичными доказательствами.

\paragraph{Допущение об исследующих стартах.}
\index{exploring starts}
Сразу возникает ключевая трудность: если стратегия $\pi$ детерминирована, она никогда не выберет определённые действия в определённых состояниях, поэтому пары $(s, a)$ для этих действий никогда не посещаются и $Q^\pi(s, a)$ никогда не оценивается.  Без оценок для всех пар нельзя улучшить стратегию, выбирая жадное действие.

Допущение об \emph{исследующих стартах} разрешает эту проблему, требуя, чтобы каждая пара состояние--действие имела положительную вероятность быть стартовой парой эпизода:
\[
  \Prob\!\bigl(S_0 = s,\; A_0 = a\bigr) > 0 \quad \forall\, s \in \cS,\; a \in \cA.
\]
При исследующих стартах каждая пара $(s, a)$ будет в конечном счёте посещаться бесконечно часто (поскольку эпизоды генерируются многократно), обеспечивая состоятельную оценку $Q^\pi$ для всех пар.  Раздел~\ref{sec:mc-es} разрабатывает полный алгоритм управления на основе этого допущения; раздел~\ref{sec:mc-on-policy} затем снимает его.


% ===========================================================
\section{Управление методом Монте-Карло с исследующими стартами}
\label{sec:mc-es}
\index{Monte Carlo control!exploring starts}

Фреймворк обобщённой итерации по стратегии (GPI) из главы~\ref{ch:dp} чередует \emph{оценивание стратегии} (оценку $Q^\pi$) и \emph{улучшение стратегии} (выбор жадной стратегии по текущей $Q$).  Управление методом Монте-Карло реализует этот фреймворк, используя выборочные эпизоды вместо обходов уравнений Беллмана.

\begin{algorithm}[ht]
\caption{Управление методом МК с исследующими стартами (MC-ES)}
\label{alg:mc-es}
\KwIn{Дисконт $\gamma$}
Инициализировать $Q(s,a)$ произвольно, $\pi(s) \in \cA$ произвольно, $\mathit{Count}(s,a) \leftarrow 0$ для всех $s, a$\;
\For{каждый эпизод}{
  Выбрать $S_0 \sim \mbox{Равном.}(\cS)$, $A_0 \sim \mbox{Равном.}(\cA)$ \upshape(исследующие старты)\;
  Сгенерировать эпизод $S_0, A_0, R_1, S_1, A_1, R_2, \dots, S_\tau$ следуя $\pi$ после шага 0\;
  $G \leftarrow 0$\;
  \For{$t = \tau-1, \tau-2, \dots, 0$ \upshape(в обратном порядке)}{
    $G \leftarrow \gamma G + R_{t+1}$\;
    \If{$(S_t, A_t)$ не встречается в $(S_0,A_0),\dots,(S_{t-1},A_{t-1})$}{
      $\mathit{Count}(S_t, A_t) \leftarrow \mathit{Count}(S_t, A_t) + 1$\;
      $Q(S_t, A_t) \leftarrow Q(S_t, A_t) + \frac{1}{\mathit{Count}(S_t,A_t)}\bigl(G - Q(S_t, A_t)\bigr)$\;
      $\pi(S_t) \leftarrow \argmax_{a} Q(S_t, a)$\;
    }
  }
}
\end{algorithm}

\paragraph{Аргумент сходимости.}
Алгоритм выполняет один шаг улучшения стратегии (строка 11) после каждого обновления отдачи.  Формально, сходимость к оптимальной стратегии $\pi^*$ следует из аргумента о неподвижной точке: если алгоритм сходится, то в точке сходимости $Q \approx Q^{\pi}$ (оценивание состоятельно по теореме~\ref{thm:fv-unbiased}) и $\pi$ жадна по отношению к $Q^{\pi}$.  Стратегия, одновременно жадная по отношению к собственной функции ценности, должна быть оптимальной (см. теорему об улучшении стратегии, глава~\ref{ch:dp}).

\paragraph{Почему исследующие старты непрактичны.}  Допущение о том, что каждый эпизод может начинаться с произвольной пары $(s, a)$, искусственно в большинстве реальных задач.  Например, нельзя начать шахматную партию с королём в произвольной позиции или инициализировать робота в произвольном положении суставов, гарантируя при этом безопасность.  Требование исследования должно быть встроено в саму стратегию\,---\,это тема следующих двух разделов.


% ===========================================================
\section{Управление методом Монте-Карло на текущей стратегии}
\label{sec:mc-on-policy}
\index{Monte Carlo control!on-policy}

Управление МК на текущей стратегии отказывается от исследующих стартов и вместо этого обеспечивает попытку каждого действия в каждом состоянии, используя \emph{стохастическую} стратегию, которая всегда назначает некоторую минимальную вероятность каждому действию.

\begin{definition}[$\varepsilon$-мягкая стратегия]
\label{def:eps-soft}
\index{$\varepsilon$-soft policy}
Стратегия $\pi$ называется \textbf{$\varepsilon$-мягкой}, если для всех $s \in \cS$ и всех $a \in \cA(s)$:
\begin{equation}
  \label{eq:eps-soft}
  \pi(a \mid s) \ge \frac{\varepsilon}{|\cA(s)|}, \qquad \varepsilon > 0.
\end{equation}
\end{definition}

Среди всех $\varepsilon$-мягких стратегий естественным улучшением данной $\varepsilon$-мягкой стратегии $\pi$ является $\varepsilon$-жадная стратегия по отношению к текущей оценке $Q$:
\begin{equation}
  \label{eq:eps-greedy-policy}
  \pi'(a \mid s) = \begin{cases}
    1 - \varepsilon + \dfrac{\varepsilon}{|\cA(s)|}, & a = \argmax_{b} Q(s, b), \\[6pt]
    \dfrac{\varepsilon}{|\cA(s)|}, & \mbox{иначе.}
  \end{cases}
\end{equation}

\subsection*{Теорема об улучшении $\varepsilon$-мягкой стратегии}

\begin{theorem}[Улучшение $\varepsilon$-мягкой стратегии]
\label{thm:eps-soft-improvement}
\index{$\varepsilon$-soft policy improvement theorem}
Пусть $\pi$ --- $\varepsilon$-мягкая стратегия и $\pi'$ --- $\varepsilon$-жадная стратегия, определённая формулой~\eqref{eq:eps-greedy-policy} по отношению к $Q^\pi$.  Тогда $v^{\pi'}(s) \ge v^\pi(s)$ для всех $s \in \cS$.
\end{theorem}

\begin{proof}
Зафиксируем $s \in \cS$ и пусть $m = |\cA(s)|$.  Обозначим $M = \max_{a \in \cA(s)} Q^\pi(s, a)$.

Поскольку $\pi$ является $\varepsilon$-мягкой, можно разложить её как
\[
  \pi(a \mid s) = \frac{\varepsilon}{m} + \delta_a, \quad \delta_a \ge 0, \quad \sum_a \delta_a = 1 - \varepsilon.
\]
Тогда, используя сокращённое обозначение $v^\pi(s) = Q^\pi(s, \pi) := \sum_a \pi(a \mid s) Q^\pi(s, a)$:
\begin{align}
  v^\pi(s)
  &= \sum_{a} \pi(a \mid s)\,Q^\pi(s, a) \notag \\
  &= \frac{\varepsilon}{m} \sum_a Q^\pi(s, a) + \sum_a \delta_a\,Q^\pi(s, a) \notag \\
  &\le \frac{\varepsilon}{m} \sum_a Q^\pi(s, a) + (1 - \varepsilon)\,M. \label{eq:eps-soft-ineq}
\end{align}
Заметим, что правая часть~\eqref{eq:eps-soft-ineq} есть в точности
\[
  Q^\pi(s, \pi') = \sum_a \pi'(a \mid s)\,Q^\pi(s, a),
\]
поскольку $\pi'$ назначает вес $1 - \varepsilon + \varepsilon/m$ максимизирующему действию и $\varepsilon/m$ всем остальным.  Таким образом,
\[
  Q^\pi(s, \pi') \ge v^\pi(s) = Q^\pi(s, \pi).
\]
Это в точности условие теоремы об улучшении стратегии (глава~\ref{ch:dp}), которая даёт $v^{\pi'}(s) \ge v^\pi(s)$.
\end{proof}

\begin{remark}
Неравенство превращается в равенство\,---\,и тем самым $\pi = \pi'$\,---\,только когда $Q^\pi(s, a)$ одинаково для всех $a$ в каждом состоянии, то есть когда $\pi$ уже оптимальна среди всех $\varepsilon$-мягких стратегий.  Этот оптимум --- \emph{$\varepsilon$-оптимальная}\index{$\varepsilon$-optimal policy} стратегия: наилучшая достижимая при условии сохранения не менее $\varepsilon/|\cA(s)|$ вероятности на каждое действие.  В общем случае это не глобально оптимальная детерминированная стратегия $\pi^*$.  Разрыв уменьшается при $\varepsilon \to 0$.
\end{remark}

\begin{algorithm}[ht]
\caption{Управление МК по первому посещению на текущей стратегии ($\varepsilon$-жадная)}
\label{alg:on-policy-mc}
\KwIn{Малое $\varepsilon > 0$, дисконт $\gamma$}
Инициализировать $Q(s,a) \leftarrow 0$, $\mathit{Count}(s,a) \leftarrow 0$ для всех $s, a$\;
$\pi \leftarrow \varepsilon$-жадная по $Q$\;
\For{каждый эпизод}{
  Сгенерировать эпизод $S_0, A_0, R_1, \dots, S_\tau$ по стратегии $\pi$\;
  $G \leftarrow 0$\;
  \For{$t = \tau-1, \tau-2, \dots, 0$}{
    $G \leftarrow \gamma G + R_{t+1}$\;
    \If{$(S_t, A_t)$ первое посещение в эпизоде}{
      $\mathit{Count}(S_t, A_t) \leftarrow \mathit{Count}(S_t, A_t) + 1$\;
      $Q(S_t, A_t) \leftarrow Q(S_t, A_t) + \frac{1}{\mathit{Count}(S_t,A_t)}\bigl(G - Q(S_t, A_t)\bigr)$\;
    }
  }
  Обновить $\pi$ до $\varepsilon$-жадной по текущей $Q$\;
}
\end{algorithm}

Свойство $\varepsilon$-мягкости гарантирует, что каждая пара $(s, a)$ посещается бесконечно часто, поэтому оценки $Q(s, a)$ сходятся к $Q^\pi$ по теореме~\ref{thm:fv-unbiased}, и многократное применение теоремы~\ref{thm:eps-soft-improvement} толкает стратегию к $\varepsilon$-оптимальной.

\paragraph{Сходимость к $\varepsilon$-оптимальности.}  Управление МК на текущей стратегии сходится к наилучшей возможной стратегии среди всех $\varepsilon$-мягких стратегий.  Это слабее, чем сходимость к глобально оптимальной $\pi^*$.  Чтобы получить $\pi^*$, используют либо исследующие старты (раздел~\ref{sec:mc-es}), либо, что более практично, подход на другой стратегии из раздела~\ref{sec:mc-importance-sampling}, который позволяет целевой стратегии быть детерминированной, пока поведенческая стратегия занимается исследованием.


% ===========================================================
\section{Метод Монте-Карло на другой стратегии с взвешиванием по важности}
\label{sec:mc-importance-sampling}
\index{off-policy Monte Carlo}\index{importance sampling}

\subsection*{Проблема обучения на другой стратегии}

При обучении на другой стратегии данные собираются \emph{поведенческой стратегией}\index{behavior policy} $\mu$ (которая может свободно исследовать), тогда как мы хотим оценить или улучшить \emph{целевую стратегию}\index{target policy} $\pi$ (которая может быть детерминированной и жадной).  Фундаментальное препятствие состоит в том, что распределение эпизодов, порождаемое $\mu$, отличается от распределения, порождаемого $\pi$.  Наивное усреднение отдач давало бы оценки $Q^\mu$, а не $Q^\pi$.

\paragraph{Допущение о покрытии.}
\index{coverage assumption}
Чтобы оценивание на другой стратегии было корректным, требуется, чтобы $\pi$ не могла выбрать ни одно действие, которое $\mu$ никогда не выбирает:
\begin{equation}
  \label{eq:coverage}
  \pi(a \mid s) > 0 \implies \mu(a \mid s) > 0, \quad \forall\, s \in \cS,\; a \in \cA.
\end{equation}
Без~\eqref{eq:coverage} некоторые траектории, требуемые $\pi$, имели бы нулевую вероятность при $\mu$ и никогда не появлялись бы в данных.

\subsection*{Отношение важности выборок}

Ключевым инструментом является \emph{отношение важности выборок}\index{importance sampling ratio}, которое перевзвешивает вероятность наблюдения конкретной траектории при $\mu$ к вероятности наблюдения её при $\pi$.  Для отрезка траектории $(S_t, A_t, \dots, S_T)$, сгенерированного с помощью $\mu$, отношение вероятности траектории при $\pi$ к вероятности при $\mu$ равно
\begin{equation}
  \label{eq:is-ratio}
  \rho_{t:T-1} = \prod_{k=t}^{T-1} \frac{\pi(A_k \mid S_k)}{\mu(A_k \mid S_k)}.
\end{equation}
Вероятности переходов среды $p(s', r \mid s, a)$ сокращаются в этом отношении, поскольку они одинаковы для обеих стратегий.

\subsection*{Обычное взвешивание по важности (OIS)}

\begin{definition}[Оценщик обычного взвешивания по важности]
\label{def:ois}
Для $n$ эпизодов с отдачами $G_0^{(1)}, \dots, G_0^{(n)}$ и соответствующими весами ВВ $\rho^{(1)}, \dots, \rho^{(n)}$, \textbf{оценщик обычного взвешивания по важности} (OIS) для $Q^\pi(s, a)$ равен
\begin{equation}
  \label{eq:ois-estimator}
  \widehat{Q}_n^{\,\mathrm{OIS}}(s, a) = \frac{1}{n} \sum_{i=1}^{n} \rho^{(i)} G_0^{(i)}.
\end{equation}
\end{definition}

\begin{theorem}[OIS несмещён]
\label{thm:ois-unbiased}
При допущении о покрытии~\eqref{eq:coverage},
\[
  \E_\mu\!\left[\rho_{0:T-1}\,G_0 \mid S_0 = s,\; A_0 = a\right] = Q^\pi(s, a).
\]
\end{theorem}

\begin{proof}
Обозначим хвостовую траекторию через $h = (s_1, a_1, r_2, \dots, s_\tau)$.  Тогда
\[
  Q^\pi(s, a) = \sum_{h} \Prob_\pi(h \mid s, a)\,G(h),
\]
где $G(h) = \sum_{k} \gamma^k r_{k+1}$.  При стратегии $\mu$:
\[
  \Prob_\mu(h \mid s, a) = \prod_{k=1}^{\tau-1} \mu(a_k \mid s_k)\,p(s_{k+1}, r_{k+1} \mid s_k, a_k).
\]
По определению $\rho$:
\[
  \rho(h)\,\Prob_\mu(h \mid s, a)
  = \prod_{k=1}^{\tau-1} \pi(a_k \mid s_k)\,p(s_{k+1}, r_{k+1} \mid s_k, a_k)
  = \Prob_\pi(h \mid s, a).
\]
Следовательно,
\begin{align*}
  \E_\mu[\rho\,G_0 \mid S_0 = s, A_0 = a]
  &= \sum_h \rho(h)\,G(h)\,\Prob_\mu(h \mid s, a) \\
  &= \sum_h G(h)\,\Prob_\pi(h \mid s, a)
  = Q^\pi(s, a). \qedhere
\end{align*}
\end{proof}

\subsection*{Взвешенное взвешивание по важности (WIS)}

\begin{definition}[Оценщик взвешенного взвешивания по важности]
\label{def:wis}
Оценщик \textbf{взвешенного взвешивания по важности} (WIS) равен
\begin{equation}
  \label{eq:wis-estimator}
  \widehat{Q}_n^{\,\mathrm{WIS}}(s, a) = \frac{\sum_{i=1}^n \rho^{(i)}\,G_0^{(i)}}{\sum_{i=1}^n \rho^{(i)}},
\end{equation}
когда знаменатель положителен.
\end{definition}

WIS является оценщиком-отношением.  Он смещён при конечном $n$, однако асимптотически несмещён и имеет значительно меньшую дисперсию, чем OIS, в случаях, когда веса $\rho^{(i)}$ сильно варьируются.  На практике WIS обычно предпочтительнее в задачах с длинными эпизодами или близкими к детерминированной целевой стратегией, поскольку веса ВВ могут быть экспоненциально велики.

\subsection*{Инкрементная реализация WIS}

Как и в оценивании на текущей стратегии, WIS можно вычислять инкрементно, поддерживая текущую сумму $C(s,a)$ весов и текущую оценку $Q(s,a)$:
\begin{align}
  \label{eq:wis-incremental}
  C(s,a) &\leftarrow C(s,a) + \rho, \\
  Q(s,a) &\leftarrow Q(s,a) + \frac{\rho}{C(s,a)}\bigl(G - Q(s,a)\bigr).
\end{align}
Эти обновления применяются за одно обратное прохождение по каждому эпизоду, аналогично алгоритму~\ref{alg:fv-mc-prediction}.

\begin{algorithm}[ht]
\caption{МК на другой стратегии с WIS для оценивания $Q^\pi$}
\label{alg:off-policy-mc-wis}
\KwIn{Целевая стратегия $\pi$, поведенческая стратегия $\mu$, дисконт $\gamma$}
Инициализировать $Q(s,a) \leftarrow 0$, $C(s,a) \leftarrow 0$ для всех $s,a$\;
\For{каждый эпизод}{
  Сгенерировать эпизод $S_0,A_0,R_1,\dots,S_\tau$ по стратегии $\mu$\;
  $G \leftarrow 0$;\quad $\rho \leftarrow 1$\;
  \For{$t = \tau-1, \tau-2, \dots, 0$ \upshape(в обратном порядке)}{
    $G \leftarrow \gamma G + R_{t+1}$\;
    $\rho \leftarrow \rho \cdot \dfrac{\pi(A_t \mid S_t)}{\mu(A_t \mid S_t)}$\;
    $C(S_t,A_t) \leftarrow C(S_t,A_t) + \rho$\;
    $Q(S_t,A_t) \leftarrow Q(S_t,A_t) + \dfrac{\rho}{C(S_t,A_t)}\bigl(G - Q(S_t,A_t)\bigr)$\;
    \If{$\pi(A_t \mid S_t) = 0$}{\textbf{break}\;}
  }
}
\Return $Q$\;
\end{algorithm}

Строка ``break'' использует тот факт, что если $\pi(A_t \mid S_t) = 0$, то $\rho = 0$ для всех более ранних шагов --- эти шаги не вносят вклада в оценку $Q^\pi$.

% ===========================================================
\section{Управление методом Монте-Карло на другой стратегии}
\label{sec:mc-off-policy-control}
\index{Monte Carlo control!off-policy}

Объединяя WIS-оценивание с GPI, получаем полный алгоритм управления на другой стратегии.  Целевая стратегия $\pi$ является жадной по $Q$, а поведенческая стратегия $\mu$ --- $\varepsilon$-жадной (или любой другой стратегией, удовлетворяющей допущению о покрытии).

\begin{algorithm}[ht]
\caption{Управление МК на другой стратегии (WIS)}
\label{alg:off-policy-mc-control}
\KwIn{Дисконт $\gamma$, параметр исследования $\varepsilon$}
Инициализировать $Q(s,a)$ произвольно, $C(s,a) \leftarrow 0$, $\pi(s) \leftarrow \argmax_a Q(s,a)$ для всех $s,a$\;
\For{каждый эпизод}{
  $\mu \leftarrow$ любая стратегия, покрывающая $\pi$ (например, $\varepsilon$-жадная по $Q$)\;
  Сгенерировать эпизод $S_0,A_0,R_1,\dots,S_\tau$ по $\mu$\;
  $G \leftarrow 0$;\quad $\rho \leftarrow 1$\;
  \For{$t = \tau-1, \tau-2, \dots, 0$}{
    $G \leftarrow \gamma G + R_{t+1}$\;
    $\rho \leftarrow \rho \cdot \frac{\pi(A_t \mid S_t)}{\mu(A_t \mid S_t)}$\;
    $C(S_t,A_t) \leftarrow C(S_t,A_t) + \rho$\;
    $Q(S_t,A_t) \leftarrow Q(S_t,A_t) + \frac{\rho}{C(S_t,A_t)}\bigl(G - Q(S_t,A_t)\bigr)$\;
    $\pi(S_t) \leftarrow \argmax_a Q(S_t,a)$\;
    \If{$A_t \ne \pi(S_t)$}{\textbf{break}\;}
  }
}
\end{algorithm}

Условие ``break'' в строке 12 использует тот факт, что жадная целевая стратегия $\pi$ будет назначать нулевую вероятность любому действию, не являющемуся жадным; таким образом, $\rho$ становится равным нулю, и более ранние шаги не вносят вклада.

% ===========================================================
\section{Сводка и перспективы}
\label{sec:mc-summary}

Методы Монте-Карло представляют собой первый большой класс безмодельных алгоритмов обучения с подкреплением.  Их ключевые свойства таковы:
\begin{itemize}[leftmargin=*]
  \item \textbf{Несмещённость.}  Как оценщик по первому посещению (теорема~\ref{thm:fv-unbiased}), так и OIS (теорема~\ref{thm:ois-unbiased}) несмещённы; WIS асимптотически несмещён.
  \item \textbf{Большая дисперсия.}  Поскольку целевой является полная отдача $G_t$, дисперсия растёт с горизонтом.  Это главная слабость МК по сравнению с методами TD.
  \item \textbf{Нет бутстрэппинга.}  Оценки МК независимы от текущего приближения функции ценности.  Это устраняет смещение за счёт бутстрэппинга, но требует ждать конца эпизода.
  \item \textbf{Только эпизодические задачи.}  МК-методы требуют терминальных состояний.  Для продолжающихся задач необходимы методы TD.
\end{itemize}

В главе~\ref{ch:td} мы увидим, как методы временной разности устраняют последнее ограничение, жертвуя несмещённостью ради возможности обновлений после каждого отдельного перехода.

% ===========================================================
\section*{Упражнения}

\begin{enumerate}
  \item \textbf{(МК по первому посещению против каждого посещения.)}  Приведите пример МПР и стратегии, при которых дисперсия оценщика по каждому посещению строго меньше дисперсии оценщика по первому посещению для некоторого состояния $s$.  Объясните интуицию.

  \item \textbf{(Смещение WIS.)}  Для $K=2$ независимых эпизодов с отдачами $G^{(1)}, G^{(2)}$ и весами ВВ $\rho^{(1)}, \rho^{(2)}$ вычислите $\E[\widehat{Q}_2^{\mathrm{WIS}}]$ явно и покажите, что он смещён при $Q^\pi = c$ (постоянном).  Какова скорость убывания смещения при $n \to \infty$?

  \item \textbf{(Исследующие старты.)}  Объясните, почему MC-ES (алгоритм~\ref{alg:mc-es}) не гарантирует исследование всех пар $(s,a)$ в первом эпизоде, но гарантирует исследование каждой пары бесконечно часто в пределе.

  \item \textbf{(Взрыв дисперсии OIS.)}  Рассмотрим МПР, в котором $\pi$ является детерминированной и $\mu$ --- равномерной по $|\cA|$ действиям.  Покажите, что для эпизода длины $\tau$ вес ВВ удовлетворяет $\rho = |\cA|^\tau$ почти наверное.  Какой вывод следует отсюда для практического применения OIS при длинных эпизодах?

  \item \textbf{(Инкрементный WIS.)}  Докажите, что обновления~\eqref{eq:wis-incremental} алгебраически эквивалентны вычислению $\widehat{Q}_n^{\mathrm{WIS}}$ по определению~\eqref{eq:wis-estimator} после $n$ обновлений.
\end{enumerate}
